\section{Discussion}
\label{sec:discussion}

\subsection{What the theory explains that predecessors do not}

The theory provides unified explanations for phenomena that existing frameworks address only partially or not at all.

\textbf{Why style predicts more than material.} Hatchuel and Weil's C-K theory does not address morphological variation directly; Stiny's shape grammars generate forms but do not predict which forms are realized. The fitness landscape framework explains the empirical finding straightforwardly: style is a proxy for the full pressure set (minus the class-constant affordance), while material is a single pressure. The aggregate variable necessarily carries more information about form variation than any individual component.

\textbf{Why styles are gradients, not clusters.} If styles were imposed by external categorical forces (aesthetic movements, Zeitgeist), they would produce discrete clusters in morphospace. The negative silhouette scores ($-0.338$) confirm that styles are gradient phenomena: continuous overlapping distributions produced by agents navigating a shared but shifting landscape. Style is the statistical trace of navigation, not a cause of it \parencite{Kubler1962}.

\textbf{Why form history lags material history.} Steel existed for centuries before the tube chair. The theory explains this latency through the distinction between material affordance (the operations a material \textit{permits}) and the operation itself (the technique that \textit{activates} the affordance). Material availability is necessary but not sufficient; the operation must enter the knowledge space $K$ before the morphospace can expand into the newly accessible region.

\textbf{Why ``form follows function'' is empirically false within a class.} The affordance structure defines the class and is constant within it (Proposition 1). By definition, it carries zero information about variation within the class. Any variation must be driven by pressures other than the class-defining affordance. The empirical data confirm this: the class-constant variable (use/function) is informationally empty within the class, and the aggregate of non-functional pressures (style) dominates.

\subsection{Limitations}

The theory has been tested against a single artifact class (the European chair). Generalization to other classes; architecture, graphic design, software interfaces, biological design; is theorized but not yet demonstrated. Cross-domain replication is the most immediate research priority (Section~\ref{sec:agenda}).

The corpus is geographically limited to European collections. Non-European design traditions are not represented. If the theory's landscape attractors reflect universal physical and ergonomic constraints, non-European chairs should converge on similar morphospace structures. If they do not, the theory requires revision to accommodate culturally specific landscape features.

The three-dimensional mesh models are computationally generated from catalog images, not directly scanned from physical objects. While the geometric features extracted from these models produce consistent results across 84 cross-validation tests, direct 3D scanning would strengthen the empirical foundation.

The theory does not yet specify the functional form of the fitness landscape. The weights $w_i(t)$ in the landscape equation $L(c,t) = \sum w_i(t) \cdot p_i$ are treated as time-varying parameters, not as functions with specified dynamics. A fully predictive theory would require modeling the temporal evolution of these weights; this remains an open problem.

\subsection{Relationship to existing design theories}

The theory is \textit{stronger} than C-K theory alone: it adds morphospace geometry and landscape dynamics to the epistemic logic, providing $C$ with multi-dimensional structure and $K$ with empirically measurable content.

It is \textit{more general} than shape grammars alone: it adds the agent, the selection pressure, and the landscape to the syntactic apparatus, explaining not only what operations are possible but why certain operations are performed.

It is \textit{more specific} than Levin's multi-scale competence framework alone: it provides a design-domain formalization (morphospace, affordance, grammar, style) that makes the general framework empirically operational for artifacts.

The theory does not \textit{replace} these frameworks; it \textit{integrates} them. Each tradition contributes a necessary layer: Stiny's algebra provides the syntactic substrate, Hatchuel and Weil's operators provide the epistemic dynamics, and Levin's competence hierarchy provides the causal mechanism. The core equation $\mathrm{Form} = SG(\nabla_{C(A)} L(c,t))$ is the point of integration.
